\chapter{METODOLOGI}

\section{Struktur Tabel}
Dalam percobaan ini, digunakan contoh dua tabel: tabel operasional dan tabel audit.

\begin{table}[H]
    \centering
    \caption{Skema Tabel Transaksi}
    \label{tab:skema}
    \begin{tabularx}{\textwidth}{@{} l l X @{}}
        \toprule
        \textbf{Kolom} & \textbf{Tipe Data} & \textbf{Keterangan} \\
        \midrule
        \texttt{id} & \texttt{INT} & \textit{Primary Key} \\
        \texttt{nominal} & \texttt{DECIMAL} & Nilai transaksi masuk / keluar \\
        \texttt{keterangan} & \texttt{VARCHAR} & Identitas atau detil ringkas laporan transaksi \\
        \texttt{user\_id} & \texttt{INT} & Mengacu ke aktor aplikasi (\textit{Foreign Key}) \\
        \bottomrule
    \end{tabularx}
\end{table}

Berdasarkan Tabel~\ref{tab:skema}, setiap baris memuat keterangan dan nilai uang untuk kebutuhan transaksi.

\section{Pembuatan Prosedur Dasar}
Kode berikut akan diujicobakan pada ekosistem basis data.

\begin{lstlisting}[language=SQL, caption={Contoh Query trigger audit}]
CREATE TRIGGER trg_setelah_update
AFTER UPDATE ON transaksi
FOR EACH ROW
BEGIN
    INSERT INTO transaksi_audit (id_trans, aksi, waktu)
    VALUES (OLD.id, 'UBAH', NOW());
END;
\end{lstlisting}

\section{Pembuatan Prosedur Lanjutan Aplikasi}

Pada aplikasi, panggilan kueri diproses dalam blok \textcite{masri2021}, agar dapat dievaluasi apabila status rekam gagal tereksekusi.

\begin{lstlisting}[language=TypeScript, caption={Contoh integrasi kueri pada aplikasi TypeScript}]
async function updateTransaction(id: number, money: number) {
  const db = await connect();
  try {
    await db.execute(`UPDATE transaksi SET nominal = ? WHERE id = ?`, [money, id]);
    // Triggers will automatically act on DB
  } catch (error) {
    console.error("Gagal melakukan ubah.", error);
  }
}
\end{lstlisting}
